% Aula 08 --- doença rara: o classificador de Bayes é o trivial, e os custos mudam isso.
Cecília é médica e coordena o rastreamento de uma doença que atinge $1\%$ da população:
$\Prob(Y=1)=0{,}01$, com $Y=1$ para doente. O exame disponível tem resultado $X\in\{0,1\}$ e,
seguindo o resultado (isto é, declarando doente quem tem $X=1$), detecta $90\%$ dos doentes
e dá alarme falso para $10\%$ dos saudáveis:
$\Prob(X=1\mid Y=1)=0{,}9$ e $\Prob(X=1\mid Y=0)=0{,}1$. Seja
$\eta(x)=\Prob(Y=1\mid X=x)$.
\begin{itens}
  \item \pts{1,2} Calcule $\eta(1)$ e $\eta(0)$.
  \item \pts{1,2} Determine o classificador de Bayes, o que minimiza a probabilidade de erro.
        Compare a probabilidade de erro dele com a do classificador que segue o exame. Um
        colega conclui que o exame não serve para nada. O que há de errado nessa conclusão?
  \item \pts{1,6} Sejam $\ell_0$ o custo de deixar passar um doente e $\ell_1$ o de um alarme
        falso, e $r=\ell_0/\ell_1$. Lembre que o classificador que declara doente quando
        $\eta(x)>\ell_1/(\ell_0+\ell_1)$ minimiza o custo esperado. Mostre que esse
        classificador segue o exame, isto é, declara doente exatamente quem tem $X=1$, se e
        somente se $11<r\le891$, e diga o que ele faz fora dessa faixa. Para $\ell_0=20$ e
        $\ell_1=1$, calcule o custo esperado de declarar todos saudáveis, de seguir o exame e
        de declarar todos doentes, numa população de $10\,000$ pessoas.
\end{itens}

\begin{solucao}
\textbf{(a)} Pelo Teorema de Bayes,
\[
  \eta(1)=\frac{0{,}9\cdot0{,}01}{0{,}9\cdot0{,}01+0{,}1\cdot0{,}99}
        =\frac{0{,}009}{0{,}108}=\frac1{12}\approx0{,}083,
\]
\[
  \eta(0)=\frac{0{,}1\cdot0{,}01}{0{,}1\cdot0{,}01+0{,}9\cdot0{,}99}
        =\frac{0{,}001}{0{,}892}=\frac1{892}\approx0{,}0011 .
\]
Um resultado positivo multiplica a probabilidade de doença por cerca de $8$, de $1\%$ para
$8{,}3\%$, e um negativo a divide por cerca de $9$; ainda assim, mesmo com o exame positivo,
a doença continua improvável, porque ela é rara: entre os positivos, os $9$ doentes em cada
mil pessoas são superados pelos $99$ saudáveis com alarme falso.

\smallskip
\textbf{(b)} O classificador de Bayes declara doente quando $\eta(x)>1/2$. Como
$\eta(0)$ e $\eta(1)$ estão abaixo de $1/2$, ele é $g\equiv0$: declara todos saudáveis,
qualquer que seja o resultado do exame. A probabilidade de erro dele é
$\Prob(Y=1)=0{,}01$. Seguir o exame erra nos doentes com $X=0$ e nos saudáveis com $X=1$:
\[
  \Prob(X\ne Y)=0{,}01\cdot0{,}1+0{,}99\cdot0{,}1=0{,}001+0{,}099=0{,}1 ,
\]
dez vezes mais. A conclusão do colega confunde a probabilidade de erro com a utilidade. O
classificador de Bayes minimiza a probabilidade de erro, que dá o mesmo peso a um doente não
detectado e a um alarme falso; com uma doença rara, a maneira mais barata de errar pouco é
nunca acusar ninguém, ao preço de não detectar doente nenhum. O exame detecta $90\%$ dos
doentes: o que falta saber é quanto vale cada um deles comparado a um alarme falso.

\smallskip
\textbf{(c)} O corte é $\ell_1/(\ell_0+\ell_1)=1/(1+r)$, e o classificador compara $\eta(x)$
com ele em cada valor de $x$ separadamente. Ele segue o exame, isto é, declara doente em
$x=1$ e saudável em $x=0$, quando
\[
  \eta(1)>\frac1{1+r}\iff1+r>12\iff r>11
  \qquad\text{e}\qquad
  \eta(0)\le\frac1{1+r}\iff r\le891 .
\]
Com $r\le11$, ele declara todos saudáveis; com $r>891$, declara todos doentes, porque,
quando um doente não detectado custa mais que $891$ alarmes falsos, até a chance de $1/892$
que resta depois de um exame negativo justifica declarar doente. Nos extremos da faixa há
empate: com $r=11$, declarar doente ou saudável custa o mesmo em $x=1$, e, com $r=891$, em
$x=0$. Nos dois empates o classificador declara saudável, porque só declara doente quando
$\eta(x)$ passa do corte; por isso a faixa exclui $r=11$ e inclui $r=891$. Em termos de
custo, seguir o exame é uma decisão ótima para $11\le r\le891$, e a única para
$11<r<891$: vale a pena exatamente quando deixar passar um doente custa entre $11$ e $891$
alarmes falsos.

Com $\ell_0=20$ e $\ell_1=1$, $r=20$ está na faixa, e seguir o exame é o ótimo. Numa
população de $10\,000$ pessoas há, em média, $P=100$ doentes e $N=9\,900$ saudáveis, e as
contagens esperadas saem das taxas: $\mathrm{FN}=(1-\mathrm{TPR})\,P$ e
$\mathrm{FP}=\mathrm{FPR}\cdot N$. Declarar todos saudáveis dá $\mathrm{FN}=100$ e custo
$20\cdot100=2\,000$. Seguir o exame dá $\mathrm{FN}=0{,}1\cdot100=10$ e
$\mathrm{FP}=0{,}1\cdot9\,900=990$, com custo $20\cdot10+990=1\,190$. Declarar todos doentes dá
$\mathrm{FP}=9\,900$ e custo $9\,900$. O exame economiza $810$ unidades de custo a cada dez mil
pessoas em relação ao classificador de Bayes, embora erre dez vezes mais que ele.
\end{solucao}
